\daytitle{26}{热学基础}

\begin{learningbox}{本讲义大致内容}
热学研究大量分子的集体行为。我们先从分子动理论理解温度、内能和压强，再学习理想气体实验定律，最后用热力学第一定律分析吸热、做功和内能变化。重点不是把热学公式背下来，而是建立“宏观现象由微观运动解释”的视角。
\end{learningbox}

\section*{一、学习目标}
\begin{enumerate}
  \item 理解分子动理论的三条基本观点。
  \item 区分温度、热量、内能和做功。
  \item 掌握玻意耳定律、查理定律、盖-吕萨克定律及其图像。
  \item 会用理想气体状态方程分析简单变化过程。
  \item 会用热力学第一定律判断内能变化。
\end{enumerate}

\section*{二、初高中衔接}

初中热学常把“热”与“温度升高”联系起来，但高中必须更精确：温度是状态量，热量是过程量；内能属于系统，热量和功是改变内能的两种方式。说“物体含有热量”是不严谨的，应说“物体具有内能”，或“在某过程中吸收、放出热量”。

\textbf{科学发展视角：} 热学是能源技术的基础。热机效率、空调制冷、发动机、相变储能、温室效应、低温物理和芯片散热，都要求我们理解能量如何在微观粒子运动和宏观做功之间转化。

\section*{三、核心概念}

\begin{conceptbox}{分子动理论}

分子动理论用微观运动解释宏观热现象，包含三个基本观点：
\begin{enumerate}
  \item 物体由大量分子组成；
  \item 分子永不停息地做无规则热运动；
  \item 分子间存在相互作用力。
\end{enumerate}

\begin{center}
% TikZ 图：分子动理论
% Day 26：热学基础
\begin{tikzpicture}[scale=0.9,>=Latex]
  \draw[rounded corners=3pt,thick] (0,0) rectangle (4.5,2.8);
  \foreach \x/\y/\dx/\dy in {0.7/0.8/0.55/0.35,1.3/2.0/-0.45/0.25,2.0/1.2/0.5/-0.35,3.1/2.2/-0.4/-0.3,3.7/0.75/0.45/0.4,2.9/1.35/0.35/0.45}{
    \fill[blue!55] (\x,\y) circle (2.5pt);
    \draw[->,blue!70] (\x,\y) -- ++(\dx,\dy);
  }
  \node at (2.25,-0.35) {分子不停做无规则热运动};
  \node[align=left,font=\small] at (7.0,1.45)
    {对理想气体，温度越高，\\分子平均平动动能越大。\\[3pt]
     宏观压强来自\\大量分子碰撞器壁。};
\end{tikzpicture}


\figcaption{分子热运动：宏观温度反映大量分子运动的统计特征}\label{fig:molecular-motion}
\end{center}

\textbf{关键理解：} 温度越高，分子平均动能越大；但不能说某一个分子的动能一定变大，因为分子运动是统计性的。
\end{conceptbox}

\begin{conceptbox}{温度、内能、热量}

\begin{itemize}
  \item \textbf{温度：} 描述分子平均动能大小的宏观量。
  \item \textbf{内能：} 物体内所有分子的动能和分子势能总和。
  \item \textbf{热量：} 热传递过程中转移的能量，是过程量。
  \item \textbf{做功：} 通过宏观力和位移改变系统能量的方式。
\end{itemize}

\textbf{严谨表述：} 不能说“物体含有热量”，只能说“物体吸收热量”或“放出热量”。
\end{conceptbox}

\begin{conceptbox}{气体实验定律}

一定质量理想气体在某个状态下可用 $p,V,T$ 描述。温度必须使用热力学温度：
\[
T=t+273.15.
\]

\begin{center}
% TikZ 图：气体实验定律图像
% Day 26：热学基础
\begin{tikzpicture}[scale=0.82,>=Latex]
  \begin{scope}
    \draw[->,thick] (0,0) -- (2.4,0) node[right] {$V$};
    \draw[->,thick] (0,0) -- (0,2.2) node[above] {$p$};
    \draw[blue,thick,samples=60,domain=0.35:2.25] plot(\x,{0.75/\x+0.2});
    \node at (1.2,-0.55) {等温：$pV=C$};
  \end{scope}
  \begin{scope}[xshift=3.5cm]
    \draw[->,thick] (0,0) -- (2.4,0) node[right] {$T$};
    \draw[->,thick] (0,0) -- (0,2.2) node[above] {$V$};
    \draw[red,thick] (0.25,0.25) -- (2.1,1.85);
    \node at (1.2,-0.55) {等压：$\dfrac{V}{T}=C$};
  \end{scope}
  \begin{scope}[xshift=7.0cm]
    \draw[->,thick] (0,0) -- (2.4,0) node[right] {$T$};
    \draw[->,thick] (0,0) -- (0,2.2) node[above] {$p$};
    \draw[green!55!black,thick] (0.25,0.25) -- (2.1,1.85);
    \node at (1.2,-0.55) {等容：$\dfrac{p}{T}=C$};
  \end{scope}
\end{tikzpicture}


\figcaption{气体实验定律图像：温度必须用热力学温度 $T$}\label{fig:gas-laws-graphs}
\end{center}

\[
\boxed{pV=C\quad(\text{等温})}
\]
\[
\boxed{\frac{V}{T}=C\quad(\text{等压})}
\]
\[
\boxed{\frac{p}{T}=C\quad(\text{等容})}
\]

综合起来，一定质量理想气体满足
\[
\boxed{\frac{pV}{T}=C}.
\]
\end{conceptbox}

\begin{conceptbox}{热力学第一定律}

取“系统吸热为正，系统对外做功为正”的约定：
\[
\boxed{\Delta U=Q-W}.
\]

\begin{center}
% TikZ 图：热力学第一定律
% Day 26：热学基础
\begin{tikzpicture}[scale=0.9,>=Latex]
  \draw[thick] (0,0) rectangle (3.2,2.2);
  \draw[fill=gray!20,thick] (2.25,0) rectangle (2.55,2.2);
  \foreach \x/\y in {0.5/0.5,0.9/1.6,1.4/0.9,1.8/1.8,2.0/0.45}{
    \fill[blue!55] (\x,\y) circle (2pt);
  }
  \draw[->,red,thick] (-1.0,1.1) -- (0,1.1) node[midway,above] {$Q>0$};
  \draw[->,blue,thick] (2.55,1.1) -- (3.65,1.1) node[midway,above] {$W>0$};
  \node at (1.35,-0.35) {气体};
  \node[align=left,font=\small] at (6.0,1.15)
    {$\Delta U=Q-W$\\[4pt]
     吸热增大内能；\\对外做功减小内能。};
\end{tikzpicture}


\figcaption{热力学第一定律：吸热和做功共同决定内能变化}\label{fig:first-law-piston}
\end{center}

\begin{itemize}
  \item $Q>0$：系统吸热；$Q<0$：系统放热。
  \item $W>0$：系统对外做功；$W<0$：外界对系统做功。
  \item $\Delta U>0$：内能增加；$\Delta U<0$：内能减少。
\end{itemize}
\end{conceptbox}

\section*{四、公式与适用条件}

\begin{center}
\begin{tabular}{lll}
\toprule
公式 & 物理意义 & 适用条件 \\
\midrule
$T=t+273.15$ & 摄氏温度转热力学温度 & 热学计算 \\
$pV=C$ & 等温变化 & 一定质量理想气体 \\
$\dfrac{V}{T}=C$ & 等压变化 & 一定质量理想气体 \\
$\dfrac{p}{T}=C$ & 等容变化 & 一定质量理想气体 \\
$\dfrac{pV}{T}=C$ & 状态方程 & 一定质量理想气体 \\
$\Delta U=Q-W$ & 热力学第一定律 & 吸热为正、对外做功为正 \\
$W=p\Delta V$ & 等压做功 & 压强恒定 \\
\bottomrule
\end{tabular}
\end{center}

\section*{五、典型模型}

\subsection*{模型 1：等温压缩或膨胀}

温度不变，理想气体内能不变。若气体被压缩，体积减小，压强增大：
\[
p_1V_1=p_2V_2.
\]

\subsection*{模型 2：等压升温}

压强不变时，温度升高，体积增大：
\[
\frac{V_1}{T_1}=\frac{V_2}{T_2}.
\]
注意温度必须用 K，不能直接用摄氏温度代入比例式。

\subsection*{模型 3：$p$-$V$ 图像中的功}

气体对外做功可由 $p$-$V$ 图像下方面积表示：
\[
W=\int p\,dV.
\]

\begin{center}
% TikZ 图：p-V 图像面积表示气体做功
% Day 26：热学基础
\begin{tikzpicture}[scale=0.9,>=Latex]
  \draw[->,thick] (0,0) -- (5.0,0) node[right] {$V$};
  \draw[->,thick] (0,0) -- (0,3.4) node[above] {$p$};
  \draw[blue,thick] (0.8,2.5) .. controls (2.1,2.2) and (3.2,1.5) .. (4.3,0.9);
  \fill[blue!12] (0.8,0) -- (0.8,2.5) .. controls (2.1,2.2) and (3.2,1.5) .. (4.3,0.9) -- (4.3,0) -- cycle;
  \draw[dashed] (0.8,0) -- (0.8,2.5);
  \draw[dashed] (4.3,0) -- (4.3,0.9);
  \node[below] at (0.8,0) {$V_1$};
  \node[below] at (4.3,0) {$V_2$};
  \node at (2.6,0.75) {$W=\int p\,dV$};
  \node[font=\small] at (2.6,-0.55) {$p$-$V$ 图线下方面积表示气体对外做功};
\end{tikzpicture}


\figcaption{$p$-$V$ 图像：图线下方面积表示气体对外做功}\label{fig:pv-work-area}
\end{center}

\subsection*{模型 4：绝热近似}

若过程很快或保温良好，可近似认为 $Q=0$。此时
\[
\Delta U=-W.
\]
气体快速膨胀对外做功，内能减少，温度降低；气体被快速压缩，外界对它做功，内能增加，温度升高。

\section*{六、例题精讲}

\begin{examplebox}{等温变化}
\textbf{题意分析：} 一定质量气体温度不变，初态压强 $p_1=1.0\times10^5\ \mathrm{Pa}$，体积 $V_1=2.0\ \mathrm{L}$。若体积压缩到 $V_2=0.50\ \mathrm{L}$，求末态压强。

\textbf{物理模型：} 等温变化，$pV=C$。

\textbf{解题过程：}
\[
p_1V_1=p_2V_2,
\qquad
p_2=\frac{p_1V_1}{V_2}
=\frac{1.0\times10^5\times2.0}{0.50}
=4.0\times10^5\ \mathrm{Pa}.
\]

\textbf{方法总结：} 等温压缩时体积变为原来的 $\dfrac14$，压强变为原来的 4 倍。
\end{examplebox}

\begin{examplebox}{等压升温}
\textbf{题意分析：} 一定质量气体在压强不变时，温度从 $27^\circ\mathrm{C}$ 升高到 $127^\circ\mathrm{C}$。初体积 $V_1=3.0\ \mathrm{L}$，求末体积。

\textbf{物理模型：} 等压变化，$\dfrac{V}{T}=C$。

\textbf{解题过程：}
\[
T_1=27+273=300\ \mathrm{K},\qquad T_2=127+273=400\ \mathrm{K}.
\]
\[
\frac{V_1}{T_1}=\frac{V_2}{T_2},
\qquad
V_2=V_1\frac{T_2}{T_1}=3.0\times\frac{400}{300}=4.0\ \mathrm{L}.
\]

\textbf{方法总结：} 气体定律中的温度必须用热力学温度，不能用摄氏温度直接比。
\end{examplebox}

\begin{examplebox}{热力学第一定律}
\textbf{题意分析：} 某气体吸收热量 $Q=500\ \mathrm{J}$，同时对外做功 $W=200\ \mathrm{J}$。求内能变化。

\textbf{物理模型：} 热力学第一定律。

\textbf{解题过程：}
\[
\Delta U=Q-W=500-200=300\ \mathrm{J}.
\]

\textbf{方法总结：} 吸热不一定全部转化为内能，若气体同时对外做功，内能增加量会小于吸收热量。
\end{examplebox}

\section*{七、变式训练}

\begin{enumerate}
  \item 一定质量气体等温膨胀，体积变为原来的 3 倍，压强怎样变化？
  \item 把 $-23^\circ\mathrm{C}$、$27^\circ\mathrm{C}$、$127^\circ\mathrm{C}$ 分别化为热力学温度。
  \item 气体放热 $100\ \mathrm{J}$，外界对气体做功 $300\ \mathrm{J}$，内能如何变化？
  \item 等压膨胀过程中，$p$-$V$ 图像下方面积表示什么物理量？
\end{enumerate}

\subsection*{训练提示}
\begin{enumerate}
  \item 等温时 $pV=C$，体积变为 3 倍，压强变为原来的 $\dfrac13$。
  \item 分别为 $250\ \mathrm{K}$、$300\ \mathrm{K}$、$400\ \mathrm{K}$，近似计算时常取 $273$。
  \item 放热 $Q=-100\ \mathrm{J}$，外界做功表示系统对外做功 $W=-300\ \mathrm{J}$。
  \item $p$-$V$ 图像下方面积表示气体对外做功；若体积减小，则做功为负。
\end{enumerate}

\section*{八、易错点}
\begin{mistakebox}
\begin{enumerate}
  \item \textbf{说“物体含有热量”：} 热量是过程量，不是状态量。
  \item \textbf{气体定律用摄氏温度：} 比例关系必须使用热力学温度 $T$。
  \item \textbf{混淆做功正负：} 本讲义约定系统对外做功为正，所以 $\Delta U=Q-W$。
  \item \textbf{把温度看成单个分子动能：} 温度对应大量分子的平均动能。
  \item \textbf{忽略内能还含分子势能：} 对实际物体和相变过程，分子势能变化也很重要。
\end{enumerate}
\end{mistakebox}

\section*{九、竞赛入门补充}
\begin{competitionbox}
\textbf{从微观角度看压强}

气体压强来自大量分子不断碰撞器壁。单个分子的碰撞很复杂，但大量分子的平均效果很稳定，这就是统计物理的基本思想。

在理想气体模型中，分子本身大小可忽略，分子间除碰撞外不考虑相互作用。由此可以推出
\[
pV=\frac23N\overline{E_k},
\]
其中 $N$ 是分子数，$\overline{E_k}$ 是单个分子的平均平动动能。又因为理想气体状态方程为
\[
pV=NkT,
\]
可得
\[
\overline{E_k}=\frac32kT.
\]

\textbf{意义：} 这说明温度不是模糊的冷热感觉，而是微观平均动能的量度。高中阶段不要求完整推导，但这个结果能帮助理解为什么绝对零度对应分子热运动的最低极限。
\end{competitionbox}

\section*{十、模型深化：热机、循环与不可逆性}

\subsection*{1.\quad 热机效率}

热机从高温热源吸收热量 $Q_1$，把其中一部分转化为对外做功 $W$，其余热量 $Q_2$ 放给低温热源：
\[
Q_1=W+Q_2.
\]
热机效率定义为
\[
\eta=\frac{W}{Q_1}=1-\frac{Q_2}{Q_1}.
\]

\begin{center}
% TikZ 图：热机循环和效率
% Day 26：热学基础
\begin{tikzpicture}[scale=0.9,>=Latex]
  \draw[->,thick] (0,0) -- (4.5,0) node[right] {$V$};
  \draw[->,thick] (0,0) -- (0,3.2) node[above] {$p$};
  \draw[fill=orange!16,draw=red,thick] (1.0,1.0) .. controls (1.5,2.5) and (3.0,2.6) .. (3.55,1.75)
    .. controls (3.15,0.75) and (1.55,0.55) .. (1.0,1.0);
  \draw[->,red,thick] (2.7,2.35) -- (3.05,2.1);
  \node at (2.25,1.45) {$W_{\text{循环}}$};
  \node[font=\small,align=left] at (6.55,1.65)
    {热机从高温热源吸热 $Q_1$，\\向低温热源放热 $Q_2$，\\对外输出功 $W=Q_1-Q_2$。\\[3pt]
     $\eta=\dfrac{W}{Q_1}$};
\end{tikzpicture}


\figcaption{热机循环：循环过程中系统回到初态，净功等于吸热与放热之差}\label{fig:heat-engine-cycle}
\end{center}

\textbf{重要理解：} 热机不可能把吸收的热量全部转化为功。实际热机还会因摩擦、散热和不可逆过程损失能量。

\subsection*{2.\quad 四种常见过程的能量判断}

\begin{center}
\begin{tabular}{lll}
\toprule
过程 & 特征 & 第一定律分析 \\
\midrule
等温过程 & 理想气体温度不变 & $\Delta U=0$，吸热常转化为对外做功 \\
等容过程 & 体积不变 & $W=0$，$\Delta U=Q$ \\
等压过程 & 压强不变 & $W=p\Delta V$ \\
绝热过程 & $Q=0$ & $\Delta U=-W$ \\
\bottomrule
\end{tabular}
\end{center}

\textbf{认知顺序：} 先判断过程类型，再判断 $Q$、$W$ 的正负，最后用 $\Delta U=Q-W$。不要一上来就代数运算。

\textbf{正负号约定：} 本讲义采用“系统对外做功为正”的约定，因此气体膨胀时 $W>0$，压缩时 $W<0$；吸热时 $Q>0$，放热时 $Q<0$。若别的资料采用“外界对系统做功为正”，第一定律形式会不同，但物理结论一致。做题时先看清约定，再列式。

\textbf{过程判断口诀：} 等温看温度，等容看体积，等压看压强，绝热看是否与外界交换热量。若题目没有明确过程类型，不能随意套某个特殊公式，只能从第一定律和状态方程出发分析。

\begin{examplebox}{等容加热}
\textbf{题意分析：} 一定质量理想气体被封闭在刚性容器中，吸收热量 $300\ \mathrm{J}$。求气体对外做功和内能变化。

\textbf{物理模型：} 等容过程，体积不变。

\textbf{解题过程：}
容器刚性，体积不变，所以
\[
W=0.
\]
由热力学第一定律：
\[
\Delta U=Q-W=300-0=300\ \mathrm{J}.
\]

\textbf{方法总结：} 等容过程没有体积变化，气体不对外做体积功。
\end{examplebox}

\begin{examplebox}{绝热压缩}
\textbf{题意分析：} 外界对气体做功 $200\ \mathrm{J}$，过程很快，可视为绝热。求内能变化，并判断温度如何变化。

\textbf{物理模型：} 绝热过程，$Q=0$。本讲义约定系统对外做功为正。

\textbf{解题过程：}
外界对气体做功 $200\ \mathrm{J}$，即气体对外做功
\[
W=-200\ \mathrm{J}.
\]
所以
\[
\Delta U=Q-W=0-(-200)=200\ \mathrm{J}.
\]
内能增加，理想气体温度升高。

\textbf{方法总结：} 绝热压缩会使气体升温，这也是打气筒快速打气时筒壁发热的原因之一。
\end{examplebox}

\subsection*{3.\quad 分层训练}

\begin{enumerate}
  \item \textbf{基础：} 一定质量理想气体等容放热，内能怎样变化？
  \item \textbf{提高：} 气体等压膨胀，吸热 $500\ \mathrm{J}$，对外做功 $180\ \mathrm{J}$，求内能变化。
  \item \textbf{综合：} 某热机每循环吸热 $1000\ \mathrm{J}$，放热 $650\ \mathrm{J}$，求输出功和效率。
  \item \textbf{开放思考：} 为什么冰箱能把热量从低温处搬到高温处，但必须消耗电能？
\end{enumerate}

\section*{十一、今日小结}
\begin{enumerate}
  \item 分子动理论把宏观热现象解释为大量分子的统计运动。
  \item 温度反映分子平均动能，内能是分子动能和分子势能总和。
  \item 热量和功都是过程量，都是改变内能的方式。
  \item 一定质量理想气体满足 $pV/T=C$。
  \item 热力学第一定律在本讲义约定下写作 $\Delta U=Q-W$。
\end{enumerate}
